\documentclass[11pt]{article}
\usepackage[margin=1.15in]{geometry}
\usepackage{amsmath,amsthm,amssymb}
\usepackage[hidelinks]{hyperref}

\newtheorem{proposition}{Proposition}
\newcommand{\rga}{\operatorname{RGA}^{(2)}}

\title{A double-star counterexample to the\texorpdfstring{\\}{ }TxGraffiti independence bound}
\author{}
\date{}

\begin{document}
\maketitle

For a finite simple graph $G$, let $\alpha(G)$ be its independence number and set
\[
 d_2(u)=\bigl|\{v\in V(G):\operatorname{dist}_G(u,v)\le 2\}\bigr|,
 \qquad
 \rga(G)=\sum_{uv\in E(G)}
 \frac{2\sqrt{d_2(u)d_2(v)}}{d_2(u)+d_2(v)}.
\]
Conjecture~2 in Davila~\cite[\S3.1]{davila} asserts that
$\alpha(G)\le \rga(G)$ for every nontrivial connected graph.

\begin{proposition}
There is a tree $T$ on $26$ vertices with
\[
 \rga(T)=1+\frac{12\sqrt{91}}5<24=\alpha(T).
\]
\end{proposition}

\begin{proof}
Join two vertices $a$ and $b$ by an edge and attach $12$ leaves to each.
The $24$ leaves form an independent set. Any independent set containing
one of the two centers contains at most the $12$ leaves adjacent to the
other center, besides the chosen center itself. Hence $\alpha(T)=24$.

Each center is within distance two of all $26$ vertices. A leaf is within
distance two of the $12$ leaves attached to its own center, including
itself, and of both centers, so its value of $d_2$ is $14$. The center
edge contributes $1$ to $\rga(T)$, and each of the $24$ other edges
contributes $2\sqrt{14\cdot26}/40=\sqrt{91}/10$. Therefore
\[
 \rga(T)=1+\frac{12\sqrt{91}}5<24,
\]
since $12^2\cdot91=13104<13225=115^2$.
\end{proof}

The same tree works if $d_2(u)$ counts only vertices distinct from $u$.
The center and leaf counts become $25$ and $13$, respectively. The
center edge still contributes $1$, while each leaf edge contributes
$5\sqrt{13}/19<\sqrt{91}/10$, as follows by squaring. Thus the index
under this convention is smaller still.

\begin{thebibliography}{1}
\bibitem{davila}
R. Davila, \emph{Graffiti$^3$: Compact Theory Libraries for Automated
Mathematical Discovery}, Research Square preprint, 2026.
\href{https://doi.org/10.21203/rs.3.rs-8493329/v1}{doi:10.21203/rs.3.rs-8493329/v1}.
\end{thebibliography}
\end{document}
